\chapter{三相记忆与自我的连续神经动力学实现}
\label{chap:tri_memory_neural_dynamics}

\section{从生命周期生成吸引子回到三相记忆的可实现数学结构}

上一章已经完成了对自我结构 $\Psi$ 与更慢生命周期生成吸引子 $\Omega$ 的区分，并由此建立了智能体长期身份连续性与“我正在成为什么”的方向性结构。至此，三相记忆、自我与生命周期生成方向已经在概念层形成了清晰的层级关系：
\[
h \rightarrow E \rightarrow \Theta \rightarrow \Psi \rightarrow \Omega,
\]
以及相应的时间尺度分离：
\[
\tau_h \ll \tau_E \ll \tau_\Theta \ll \tau_\Psi \ll \tau_\Omega.
\]
然而，如果这一体系只停留在概念描述层，那么它仍然不足以作为后续 AgentOS 工程实现的基础。尤其需要回答三个更严格的问题：第一，$h(t)$ 是否能够直接落到基础大语言模型的内部神经状态，而不是仅作为抽象“工作记忆”使用；第二，$E(t)$ 能否被定义为有限容量、可合并、具有时间与因果关系的向量情景图；第三，$\Theta(t)$ 与 $\Psi(t)$ 是否能够全部构造成可训练、可更新、可稳定控制的神经向量动力系统，而不是依赖符号化数据库或手工 Persona 配置。

因此，本章不继续向 $\Omega$ 方向增加新的慢变量，而是向下回到三相记忆本身，对 $h(t)$、$E(t)$、$\Theta(t)$ 与 $\Psi(t)$ 进行一次严格的数学重构。本章采用纯神经向量表述，不引入 SPC、AHC、TCE、CCM、Body Runtime 等后续控制模块，而是先建立一个最小但完整的连续多时间尺度记忆动力系统。该系统的总状态定义为
\[
\mathcal{M}(t)
=
\left(
h(t),E(t),\Theta(t),\Psi(t)
\right),
\]
其中
\[
h(t)
=
\text{当前时刻基础 LLM 全层激活状态的分布},
\]
\[
E(t)
=
\text{固定容量、具时间--因果结构的情景向量图},
\]
\[
\Theta(t)
=
\text{由长期经历慢速塑形的结构生成模型},
\]
\[
\Psi(t)
=
\text{跨越三相记忆并维持主体连续性的慢速自我结构}.
\]

本章的核心目标不是再提出一个新的“记忆模块”，而是把此前已经建立的三相记忆理论具体化为一个连续神经动力学系统。换言之，本章要完成的是从
\[
\boxed{
\text{State}
\rightarrow
\text{Trajectory}
\rightarrow
\text{Generator}
\rightarrow
\text{Self}
}
\]
到可计算形式的转换。

需要首先强调，本章中的 $h(t)$ 不等于上下文 Token 缓存。前文已经指出，Transformer 的真实工作记忆瓶颈并不只是被动存储多少 Token，而更接近前向传播过程中能够稳定维持多少有效关系结构。相干容量理论同样强调，Transformer 的真实推理边界由内部关系结构的稳定承载能力决定，而不是简单由 Context Length 决定。:contentReference[oaicite:0]{index=0} 因此，本章将 $h(t)$ 明确定义为基础模型内部表示空间上的状态分布，而不是外部文本序列本身。

\section{基础模型的双时间坐标：网络深度与生命周期时间}

设基础语言模型包含 $L$ 个 Transformer 层，第 $l$ 层隐藏维度为 $d_l$，在认知时刻 $t$，当前活动表示数量为 $N_t$。第 $l$ 层的隐藏状态矩阵写为
\[
H_l(t)
=
\begin{bmatrix}
h_{l,1}(t)\\
h_{l,2}(t)\\
\vdots\\
h_{l,N_t}(t)
\end{bmatrix}
\in
\mathbb{R}^{N_t\times d_l}.
\]
传统 Transformer 描述通常把 $H_l(t)$ 理解为一次前向传播过程中的中间隐藏状态。本章则进一步把它理解为“当前认知状态在第 $l$ 层表示空间中的一个有限采样”。

为避免把单个 Token 向量误当成整个工作记忆状态，对第 $l$ 层定义经验测度
\[
\mu_l^h(t)
=
\frac{1}{N_t}
\sum_{i=1}^{N_t}
\delta_{h_{l,i}(t)},
\]
其中 $\delta_x$ 表示以 $x$ 为支撑点的 Dirac 测度。于是完整工作记忆状态可以定义为
\[
h(t)
=
\left\{
\mu_1^h(t),
\mu_2^h(t),
\ldots,
\mu_L^h(t)
\right\}.
\]
因此，$h(t)$ 的本体不再是离散 Token 列表，而是
\[
\boxed{
h(t)
=
\text{基础模型当前全部层级表示状态的分布集合}.
}
\]

这一重新定义具有直接的理论意义。即使输入文本完全相同，只要当前召回的情景不同、自我状态不同、目标不同或者长期结构记忆不同，模型内部每层的表示分布仍可能发生显著变化。因此
\[
Input_t^{(1)}=Input_t^{(2)}
\]
并不推出
\[
h^{(1)}(t)=h^{(2)}(t).
\]
真正的当前认知状态必须定义在神经表示层，而不是输入字符串层。

为了将离散网络深度进一步转化为连续动力学，可以令
\[
s\in[0,1]
\]
表示连续“认知深度”，把离散层 $l=0,\ldots,L$ 重新参数化为
\[
x(s,t)\in\mathbb{R}^{d}.
\]
其中 $s$ 描述一次前向传播内部的表示变换，而 $t$ 描述智能体生命周期中的认知时间。于是基础模型可以写成
\[
\frac{\partial x}{\partial s}
=
F_{\Theta(t)}
\left(
x(s,t),
\mu_h(s,t),
c_E(t),
\Psi(t)
\right),
\]
其中 $F_{\Theta(t)}$ 是由当前结构记忆 $\Theta(t)$ 定义的神经向量场，$c_E(t)$ 表示从情景记忆中召回的条件表示，$\Psi(t)$ 则对当前神经动力场施加慢速自我调制。

因此整个系统天然具有两个不同时间方向：
\[
s:
\text{一次认知内部的网络深度},
\]
\[
t:
\text{跨认知事件的生命周期时间}.
\]
换句话说，本章真正研究的不是简单的 $x(t)$，而是
\[
\boxed{
x=x(s,t),
}
\]
即一个同时具有“认知深度”和“生命周期演化”两个坐标的多尺度神经系统。

\section{$h(t)$ 的连续分布动力学}

既然 $h(t)$ 被定义为神经表示状态的分布，那么对其进行连续动力学描述时，更自然的对象是概率密度
\[
\rho_h(x,s,t),
\]
并要求
\[
\int_{\mathbb{R}^{d}}
\rho_h(x,s,t)\,dx
=
1.
\]
在连续网络深度 $s$ 上，其演化可以写成一类带扩散项的连续性方程：
\[
\frac{\partial \rho_h}{\partial s}
+
\nabla_x\cdot
\left(
\rho_h F_{\Theta,\Psi,E}
\right)
=
D_h\nabla_x^2\rho_h.
\]
这里
\[
F_{\Theta,\Psi,E}
=
F_\Theta
+
F_E
+
F_\Psi
\]
表示当前神经表示所受到的三个主要驱动力：长期结构生成器 $\Theta$、情景召回 $E$ 与自我结构 $\Psi$。$D_h$ 则描述采样扰动、内部随机性、注意竞争、近似计算误差等造成的有效扩散。

在这一描述中，一次认知过程可以理解为表示分布的连续形变：
\[
\rho_h^{input}
\rightarrow
\rho_h^{semantic}
\rightarrow
\rho_h^{relational}
\rightarrow
\rho_h^{decision}.
\]
这并不意味着每一层严格对应某一个固定功能，而是强调：模型的认知状态应被理解为一个随着深度不断演化的分布流。

但是，三相记忆理论并不要求把所有 $H_l(t)$ 永久保存。若每次认知都完整存储所有层激活，则情景记忆将迅速退化成比 KV Cache 更昂贵的巨大缓存。因此必须存在一个从当前分布状态到有限情景向量的压缩算子：
\[
z_h(t)
=
\mathcal{C}_{hE}
\left[
h(t)
\right].
\]
其中
\[
z_h(t)\in\mathbb{R}^{d_E}
\]
称为当前认知事件的情景候选向量。

工程上，$\mathcal{C}_{hE}$ 可以通过 learnable query、Attention Pooling、Perceiver Resampler、Set Transformer 或多层矩统计器实现。例如，对所有层隐藏状态拼接后使用 $K$ 个可训练 query：
\[
q_k'
=
Attention
\left(
q_k,
H_1\oplus\cdots\oplus H_L
\right),
\]
最终得到
\[
z_h
=
MLP
\left(
q_1'\oplus\cdots\oplus q_K'
\right).
\]
因此
\[
\boxed{
h(t)
\rightarrow
z_h(t)
}
\]
完成的是“当前连续认知截面”向“可写入生命周期轨迹的有限事件表示”的压缩，而不是简单摘要当前文字。

\section{$E(t)$：固定容量时间--因果情景向量图}

情景记忆 $E(t)$ 不应被实现成无限增长的向量数据库。按照三相记忆原义，$E(t)$ 表示从当前时刻向过去延伸的一段具有时间连续性和事件关系的生命周期轨迹。因此，本章将其定义为一个固定容量的多关系向量图：
\[
E(t)
=
\mathcal{G}_E(t)
=
\left(
V_E(t),A_E(t)
\right),
\]
其中节点集合
\[
V_E(t)
=
\left\{
e_1,e_2,\ldots,e_M
\right\},
\]
并满足硬容量约束
\[
|V_E(t)|
\leq
M_{\max}.
\]

每一个情景节点定义为
\[
e_i
=
\left(
z_i,
t_i,
q_i,
r_i,
m_i,
n_i
\right),
\]
其中 $z_i\in\mathbb{R}^{d_E}$ 为情景向量，$t_i$ 为该情景形成的生命周期时间，$q_i$ 为当前重要度，$r_i$ 为可检索强度，$m_i$ 表示该节点经过多少次向量合并或事件压缩，$n_i$ 表示被重新访问或召回的次数。

情景图至少需要区分三种关系。第一种是时间关系：
\[
A_{ij}^{time}
=
\exp
\left(
-\frac{|t_i-t_j|}{\tau_T}
\right)
\mathbf{1}[t_i<t_j],
\]
它明确表达“事件 $i$ 先于事件 $j$”。第二种是因果关系：
\[
A_{ij}^{cause}
=
\sigma
\left[
f_c
\left(
z_i,z_j,\Delta t_{ij},a_{ij}
\right)
\right],
\]
其中 $f_c$ 是可训练的因果边预测器，$a_{ij}$ 可以包含两个事件之间的行动或外部变化。第三种是语义关系：
\[
A_{ij}^{sem}
=
\cos
\left(
z_i,z_j
\right).
\]
因此
\[
A_{ij}
=
\left(
A_{ij}^{time},
A_{ij}^{cause},
A_{ij}^{sem}
\right),
\]
形成一个 multi-relational graph。

必须强调，时间、因果和语义相似性不能压缩成一个统一标量权重。两个事件可以语义相似却没有直接因果关系，也可以因果相关却在向量空间中差异巨大。因此 $E(t)$ 的结构必须显式保持不同关系类型。

新事件 $z_h(t)$ 写入 $E$ 时，需要计算其初始重要度。先不考虑自我因素，可以定义
\[
q_i^{world}
=
\sigma
\left(
w_q^\top
\left[
z_i,g_i,\delta_i,u_i
\right]
\right),
\]
其中 $g_i$ 为当前目标相关度，$\delta_i$ 为预测误差或新颖度，$u_i$ 为后续结果或效用信号。自我 $\Psi$ 对该重要度的进一步调制将在后文定义。

情景重要度应随时间连续衰减：
\[
\frac{dq_i}{dt}
=
-\lambda_i q_i
+
\sum_k\eta_kR_{ki}(t).
\]
若没有任何再次激活，可退化为
\[
q_i(t)
=
q_i(t_i)
\exp
\left[
-\lambda_i(t-t_i)
\right].
\]
但不同事件不应采用统一衰减率，可以令
\[
\lambda_i
=
\lambda_0
\exp
\left(
-\alpha q_i-\beta n_i
\right),
\]
因此重要度越高、后续访问次数越多的情景，其有效衰减速度越慢。

当前工作记忆从 $E$ 中检索时，首先由当前状态生成 query：
\[
q_t
=
W_qz_h(t).
\]
候选节点得分不应只由向量相似度决定，而应同时考虑语义、时间、因果、重要度和自我相关性：
\[
S_i(t)
=
\alpha_s
\cos(q_t,z_i)
+
\alpha_c C_i
+
\alpha_tT_i
+
\alpha_q\log(q_i+\varepsilon)
+
\alpha_\Psi P_i^\Psi.
\]
于是
\[
p_i
=
\frac{
\exp\left(S_i/\tau_E\right)
}{
\sum_j
\exp\left(S_j/\tau_E\right)
}.
\]
在最简单实现中，可返回
\[
c_E(t)
=
\sum_{i\in TopK}
p_i z_i.
\]
但当当前任务依赖多个事件之间的时间或因果结构时，应直接返回一个局部情景子图：
\[
C_E(t)
=
\left(
[z_{i_1},\ldots,z_{i_K}],
A_E^{(K)}
\right),
\]
并由基础模型通过 cross-attention 或 graph-conditioned adapter 读取。

由于 $E$ 是固定容量结构，当
\[
|V_E|=M_{\max}
\]
而新事件继续写入时，必须执行删除、合并或结构迁移。定义情景保留分：
\[
R_i
=
\beta_q q_i
+
\beta_n\log(1+n_i)
+
\beta_c\,centrality_i
+
\beta_\Psi P_i^\Psi
-
\beta_a\,age_i
-
\beta_r\,redundancy_i.
\]
低保留分情景可以被淘汰，也可以优先与其他高相似节点进行向量级合并。

若两个情景满足
\[
\cos(z_i,z_j)
>
\theta_s
\]
且上下文相似度满足
\[
ContextSimilarity(i,j)
>
\theta_c,
\]
并且合并不会破坏关键因果分支，则可以执行
\[
e_i\oplus e_j
\rightarrow
e_{ij}.
\]
最简单的合并为
\[
z_{ij}
=
\frac{
w_i z_i+w_j z_j
}{
w_i+w_j
},
\qquad
w_i=m_iq_i.
\]
但简单向量平均会丢失局部结构，因此更合理的实现是学习一个 merge operator：
\[
z_{ij}
=
M_E
\left(
z_i,z_j,r_{ij}
\right),
\]
并定义重构损失
\[
\mathcal{L}_{merge}
=
\left\|
D_E(z_{ij},c_i)-z_i
\right\|^2
+
\left\|
D_E(z_{ij},c_j)-z_j
\right\|^2.
\]
合并同时必须进行图粗粒化。若原来存在
\[
e_a\rightarrow e_i\rightarrow e_b
\]
以及
\[
e_c\rightarrow e_j\rightarrow e_d,
\]
那么合并后的新节点 $e_{ij}$ 必须重新计算入边和出边：
\[
A_{a,ij}
=
g(A_{ai},A_{aj}),
\]
\[
A_{ij,b}
=
g(A_{ib},A_{jb}).
\]
因此，情景合并的真正数学形式是
\[
\boxed{
VectorMerge
+
GraphCoarsening,
}
\]
而不是单纯 embedding average。这一过程实际上构成 $E$ 自身的一种局部重整化机制。

\section{$\Theta(t)$：作为结构记忆的慢速生成 LLM}

三相记忆理论中最容易被误解的是 $\Theta$。如果把 $\Theta$ 实现成另一个向量数据库，那么它与 $E$ 只剩下存储周期不同，无法体现“结构记忆是生成器”的原始定义。因此，本章明确规定：
\[
\boxed{
\Theta
\neq
Database,
}
\]
而应定义成
\[
\boxed{
\Theta
=
\text{由长期情景经历慢速塑形的生成动力结构}.
}
\]

设基础预训练模型参数为
\[
W^0
=
\left\{
W_l^0
\right\}_{l=1}^{L}.
\]
若在线持续直接修改全部 $W^0$，会产生严重表示漂移、灾难性遗忘以及旧情景向量与当前表示空间失配。因此第一代工程实现中，建议保持
\[
W^0\approx const
\]
并把结构记忆实现为低秩慢速参数：
\[
W_l(t)
=
W_l^0
+
\Delta W_l^\Theta(t),
\]
其中
\[
\Delta W_l^\Theta
=
A_l^\Theta
B_l^{\Theta\top},
\]
\[
A_l^\Theta,B_l^\Theta
\in
\mathbb{R}^{d_l\times r},
\qquad
r\ll d_l.
\]
于是
\[
\Theta(t)
=
\left\{
A_l^\Theta(t),
B_l^\Theta(t)
\right\}_{l=1}^{L}
\]
真正进入模型内部动力场，并改变未来所有认知状态，而不只是提供外部检索信息。

为了表达多个稳定生成结构，还可以进一步引入结构基：
\[
\Theta
=
\left\{
\theta_1,\theta_2,\ldots,\theta_K
\right\},
\]
其中
\[
\theta_k
=
\left\{
A_{lk},B_{lk}
\right\}_{l=1}^{L}.
\]
当前工作状态通过结构路由器计算
\[
\alpha_k(t)
=
softmax
\left[
g_\Theta
\left(
z_h,c_E,\Psi
\right)
\right],
\]
从而得到
\[
\Delta W_l^\Theta(t)
=
\sum_{k=1}^{K}
\alpha_k(t)
A_{lk}B_{lk}^{\top}.
\]
因此 $\Theta$ 可以被理解为
\begin{statementbox}
StructuralMixtureOfGenerators,
\end{statementbox}
即一个可根据当前状态选择性激活的慢速结构生成器集合。

按照三相时间尺度原则，$\Theta$ 的主要学习输入不应来自单帧 $h$，而应来自经过时间过滤的 $E$：
\[
\boxed{
h
\rightarrow
E
\rightarrow
\Theta.
}
\]
如果允许
\[
h_t
\rightarrow
\Theta_{t+1}
\]
以较大增益直接发生，则瞬时噪声、一次性错误和偶然事件都可能迅速固化成长期结构。因此正常更新应写成
\[
\tau_\Theta
\frac{d\Theta}{dt}
=
-
\nabla_\Theta
\mathcal{L}_\Theta
\left(
\mathcal{B}_E,\Psi
\right),
\]
其中 $\mathcal{B}_E$ 是跨一定时间窗口采样得到的一组相关情景。

$\Theta$ 至少承担四类学习目标。第一类是未来预测：
\[
\mathcal{L}_{pred}
=
-
\log
P_\Theta
\left(
e_{t+1}
\mid
e_{\leq t}
\right).
\]
第二类是关系保持：
\[
\mathcal{L}_{rel}
=
\sum_{i,j}
\left\|
\widehat{A}_{ij}^\Theta
-
A_{ij}^E
\right\|^2.
\]
第三类是结构稳定约束：
\[
\mathcal{L}_{stability}
=
\left\|
\Theta-\Theta^{-}
\right\|_\Lambda^2.
\]
第四类是后文将定义的自我一致性损失
\[
\mathcal{L}_{self}.
\]
因此
\[
\mathcal{L}_\Theta
=
\lambda_p\mathcal{L}_{pred}
+
\lambda_r\mathcal{L}_{rel}
+
\lambda_s\mathcal{L}_{stability}
+
\lambda_\Psi\mathcal{L}_{self}.
\]

从本体上看，$E$ 与 $\Theta$ 的差异因此非常明确。$E$ 记录
\begin{statementbox}
WhatHappened?
\end{statementbox}
而 $\Theta$ 学习
\begin{statementbox}
WhatUsuallyGeneratesWhat?
\end{statementbox}
以及
\begin{statementbox}
WhatMayHappenNext?
\end{statementbox}
若 $E$ 中反复存在
\[
z_i\rightarrow z_j
\]
的结构，那么 $\Theta$ 不需要永久保存所有 $z_i,z_j$，而需要学习一个生成映射
\[
F_\Theta:
z_i
\mapsto
P(z_j\mid z_i).
\]
因此
\[
\boxed{
E=Trajectory,
\qquad
\Theta=Generator.
}
\]

更进一步，$\Theta$ 对未来的作用可以写为
\[
p_\Theta
\left(
z_{t+\Delta}
\mid
z_t,c_E,\Psi
\right),
\]
形成未来状态分布
\[
\rho_\Theta^{future}(z,\Delta t).
\]
这正是此前“结构记忆不仅表示过去，还包含未来生成结构”的神经动力学实现：
\[
\boxed{
\Theta
=
PastCompressedIntoFutureGeneratingStructure.
}
\]

\section{$\Psi(t)$：跨三相的慢速自我结构}

在上一章中，$\Psi$ 已经被定义为比 $\Theta$ 更慢、但比 $\Omega$ 更直接作用于日常认知的自我结构。本章必须进一步回答：$\Psi$ 怎样落到神经向量系统，而不退化为 system prompt、Persona JSON 或一个静态 self embedding。

按照三相记忆原始定义，$\Psi$ 在三个相态中承担不同作用：
\[
\Psi_h:
\text{$h$ 中的背景引力与状态扭曲},
\]
\[
\Psi_E:
\text{$E$ 中的经历筛选与价值权重},
\]
\[
\Psi_\Theta:
\text{$\Theta$ 中的结构空洞与组织中心}.
\]
因此，更合理的形式不是把 $\Psi$ 看成第四个普通记忆库，而是定义
\[
\Psi
=
\left(
\psi_h,\psi_E,\psi_\Theta
\right),
\]
三者共享一个慢速核心向量
\[
s_\Psi(t)
\in
\mathbb{R}^{d_\Psi}.
\]

该慢变量可以由长期 Self-related Episodic Memory 和 Structural Memory 联合编码：
\[
s_\Psi(t)
=
SelfEncoder
\left(
E^{self}(t),
\Theta^{self}(t)
\right).
\]
为了保持身份连续，其动力学必须显著慢于 $\Theta$：
\[
\tau_\Psi
\frac{ds_\Psi}{dt}
=
-s_\Psi
+
F_\Psi
\left(
s_\Psi,\bar E_\Psi,\Theta
\right),
\]
并满足
\[
\tau_\Psi\gg\tau_\Theta.
\]
离散工程实现中，也可采用
\[
s_\Psi(t+\Delta)
=
(1-\eta_\Psi)s_\Psi(t)
+
\eta_\Psi\widehat{s}_\Psi(t),
\]
其中
\[
\eta_\Psi\ll1.
\]
因此单次成功、失败、奖励或冲击不会直接重写自我。

$\Psi$ 在 $h$ 中的作用应表现为神经表示空间的背景形变，而不是一个显式 Token。例如，可由 $s_\Psi$ 为每层生成调制参数
\[
\gamma_l^\Psi,\beta_l^\Psi
=
G_l^\Psi
\left(
s_\Psi
\right),
\]
并令
\[
\widetilde{H}_l
=
\gamma_l^\Psi
\odot
H_l
+
\beta_l^\Psi.
\]
这是一种类似 FiLM 的条件调制。它意味着 $\Psi$ 不一定作为一个“关于我的句子”进入 h，而是直接改变当前认知状态在表示空间中的形成方式。

若进一步采用连续动力学语言，可定义一个自我势函数
\[
V_\Psi(x),
\]
并令当前表示流满足
\[
\frac{\partial x}{\partial s}
=
F_\Theta(x)
+
F_E(x)
-
\nabla_xV_\Psi(x).
\]
这里 $V_\Psi$ 并不是真实物理势能，而是数学上描述：符合长期主体结构的神经状态具有不同的稳定性和可达性。若 $x^\star$ 与当前 Self 高度一致，则可以令
\[
V_\Psi(x^\star)
\]
相对较低，从而使该类表示更容易成为稳定吸引区域。因此：
\[
\boxed{
\Psi_h
=
PotentialFieldOverCurrentRepresentation.
}
\]
这就是“自我在 h 中作为背景引力”的神经动力学表达。

在 $E$ 中，$\Psi$ 主要承担经历筛选器。对新事件 $e_i$ 定义 self relevance：
\[
r_i^\Psi
=
\sigma
\left[
f_\Psi
\left(
z_i,s_\Psi
\right)
\right].
\]
最终写入重要度可以定义为
\[
q_i
=
\sigma
\left(
\alpha_nNovelty_i
+
\alpha_gGoal_i
+
\alpha_uUtility_i
+
\alpha_\Psi r_i^\Psi
\right).
\]
于是两个 Agent 面对同一个外部事件 $C$ 时，可以出现
\[
q_C^{A}
\neq
q_C^{B}.
\]
这并不是因为外界事件本身不同，而是由于该事件与两个主体长期 Self 结构的关系不同。因此
\[
\boxed{
\Psi_E
=
ExperienceSelectionField.
}
\]

$\Psi$ 还应进一步调制 E 的遗忘速度：
\[
\lambda_i
=
\lambda_0
\exp
\left(
-\beta_\Psi r_i^\Psi
\right).
\]
于是高自我相关事件满足
\[
r_i^\Psi\uparrow
\Rightarrow
\lambda_i\downarrow,
\]
从而保存更久。相反，与主体无关且长期未被使用的事件衰减更快。因此 $\Psi$ 同时参与
\begin{statementbox}
WhatShouldBeRemembered?
\end{statementbox}
和
\begin{statementbox}
WhatCanBeForgotten?
\end{statementbox}
这才符合原有“自我是情景记忆筛选器”的定义。

$\Psi$ 在 $\Theta$ 中的作用则更加特殊。它不应该表现为普通世界对象节点，而应表现为长期结构的参考原点。设
\[
\Theta
=
\left\{
\theta_1,\ldots,\theta_K
\right\},
\]
不同结构基都通过
\[
c_k
=
g_\Psi
\left(
\theta_k,s_\Psi
\right)
\]
获得相对于主体的条件坐标，从而形成
\[
p_\Theta
\left(
z_{future}
\mid
z_{now},s_\Psi
\right).
\]
这里 $s_\Psi$ 不必作为一个普通 Token 出现在世界模型中，但整个生成结构都依赖于这一慢变量。因此“空洞中心”真正表达的是：Self 不一定是长期知识图中的一个普通节点，却决定许多结构如何围绕主体组织。

可以进一步把 $\Theta$ 看成一个参数流形 $\mathcal{M}_\Theta$，并定义 self-conditioned chart：
\[
\phi_\Psi:
\mathcal{M}_\Theta
\rightarrow
\mathbb{R}^{m}.
\]
各结构成分相对于自我原点重新表示为
\[
\widetilde{\theta}_i
=
\theta_i
-
C_i
\left(
s_\Psi
\right).
\]
因此
\[
\boxed{
\Psi_\Theta
=
CoordinateOriginOfPersonalStructuralMemory.
}
\]
这就是“自我在 $\Theta$ 中表现为结构空洞”的更严格数学解释：它不是某个普通结构，而是其他长期结构被个人化组织时所使用的参考中心。

为了避免长期结构在短期经验冲击下破坏身份连续，还需要让 $\Psi$ 对 $\Theta$ 的更新施加一致性正则：
\[
\mathcal{L}_{self}
=
D_\Psi
\left[
P_{\Theta'}(\cdot\mid s_\Psi),
P_{\Theta}(\cdot\mid s_\Psi)
\right].
\]
也可以采用简单函数空间约束：
\[
\mathcal{L}_{self}
=
\left\|
G_{\Theta'}(s_\Psi)
-
G_{\Theta}(s_\Psi)
\right\|^2.
\]
最终结构记忆更新为
\[
\mathcal{L}_{\Theta}^{total}
=
\mathcal{L}_{world}
+
\lambda_\Psi\mathcal{L}_{self}.
\]
因此 $\Theta$ 必须同时满足两种要求：一方面提高对世界的预测能力，另一方面不能因为局部新经验而无约束地破坏主体的长期结构连续性。

但是 $\Psi$ 本身也不能被冻结。否则 Self 将变成不可学习的静态 Persona。正常的自我更新路径应为
\[
h
\rightarrow
E
\rightarrow
\Theta
\rightarrow
\Psi,
\]
而不是
\[
h_t
\rightarrow
\Psi_{t+1}.
\]
因此可以写
\[
\Psi(t+\Delta)
=
\Psi(t)
+
\eta_\Psi
F
\left(
E_{long},
\Theta,
Outcome
\right),
\]
且要求
\[
\eta_\Psi
\ll
\eta_\Theta.
\]
只有长期重复、跨场景一致、具有高自我相关性的经验，才应逐渐推动 $\Psi$ 漂移。

\section{四层多时间尺度连续动力系统}

至此，三相记忆与自我可以统一写成一个典型的奇异摄动系统：
\[
\tau_h
\ll
\tau_E
\ll
\tau_\Theta
\ll
\tau_\Psi.
\]
定义
\[
\epsilon_E
=
\frac{\tau_h}{\tau_E},
\qquad
\epsilon_\Theta
=
\frac{\tau_E}{\tau_\Theta},
\qquad
\epsilon_\Psi
=
\frac{\tau_\Theta}{\tau_\Psi},
\]
并满足
\[
1
\gg
\epsilon_E
\gg
\epsilon_\Theta
\gg
\epsilon_\Psi
>
0.
\]

工作记忆层可以统一写成
\[
\frac{\partial\rho_h}{\partial s}
=
-
\nabla_x\cdot
\left\{
\rho_h
\left[
F_\Theta
+
F_E
-
\nabla V_\Psi
\right]
\right\}
+
D_h\nabla_x^2\rho_h.
\]

情景记忆的节点向量在写入后短时间内可近似保持稳定：
\[
\dot z_i
\approx
0,
\]
而其重要度和连接结构持续演化：
\[
\dot q_i
=
-\lambda_i(\Psi)q_i
+
R_i(t),
\]
\[
\dot A_{ij}
=
\eta_E
\left[
\widehat{A}_{ij}(h)
-
A_{ij}
\right].
\]
与此同时，E 还存在离散的 Insert、Merge 与 Evict，因此其本质是
\[
\boxed{
ContinuousWeightDynamics
+
DiscreteGraphEvents.
}
\]

结构记忆满足
\[
\tau_\Theta
\dot\Theta
=
-
\nabla_\Theta
\left[
\mathcal{L}_{pred}
+
\lambda_r\mathcal{L}_{rel}
+
\lambda_s\mathcal{L}_{stability}
+
\lambda_\Psi\mathcal{L}_{self}
\right].
\]

自我则满足
\[
\tau_\Psi
\dot s_\Psi
=
-s_\Psi
+
F_\Psi
\left(
\bar E_\Psi,\Theta,s_\Psi
\right).
\]

因此，可以将完整系统统一写成
\[
\boxed{
\mathcal{A}(t)
=
\left(
\rho_h,
\mathcal{G}_E,
\Theta,
\Psi
\right).
}
\]
其动力学为
\[
\partial_s\rho_h
=
\mathcal{F}_h
\left(
\rho_h,\mathcal{G}_E,\Theta,\Psi,x
\right),
\]
\[
\dot{\mathcal{G}}_E
=
\epsilon_E
\mathcal{F}_E
\left(
\mathcal{G}_E,\rho_h,\Psi
\right),
\]
\[
\dot\Theta
=
\epsilon_\Theta
\mathcal{F}_\Theta
\left(
\Theta,\mathcal{G}_E,\Psi
\right),
\]
\[
\dot\Psi
=
\epsilon_\Psi
\mathcal{F}_\Psi
\left(
\Psi,\mathcal{G}_E,\Theta
\right).
\]

其中，自我并不只存在于最后一条方程中，而是通过三条不同路径同时进入整个系统：
\[
\Psi
\xrightarrow{Potential}
h,
\]
\[
\Psi
\xrightarrow{Selection}
E,
\]
\[
\Psi
\xrightarrow{StructuralOrigin}
\Theta.
\]

因此，本系统不是一条
\[
h\rightarrow E\rightarrow\Theta\rightarrow\Psi
\]
的单向流水线，而是一个多时间尺度双向闭环。快状态不断向慢结构沉降，慢结构又持续塑造未来快状态。

\section{三相记忆的六个基本变换算子}

为了使上述连续动力系统工程可实现，可以进一步把三相之间的主要作用压缩成六个基本算子。

第一类是当前状态情景化算子：
\[
\mathcal{C}_{hE}:
h
\rightarrow
E.
\]
它负责把当前分布状态压缩成可写入生命周期轨迹的事件向量。

第二类是情景召回算子：
\[
\mathcal{R}_{Eh}:
E
\rightarrow
h.
\]
其目标不是恢复原始文本，而是恢复当前任务真正需要的潜在状态：
\[
\mathcal{R}_{Eh}(E)
\rightarrow
\Delta H_l.
\]
工程上可以通过
\[
H_l'
=
H_l
+
Attention
\left(
H_l,C_E,C_E
\right)
\]
把召回情景作为附加状态重新注入当前网络。

第三类是情景结构化算子：
\[
\mathcal{C}_{E\Theta}:
E
\rightarrow
\Theta.
\]
它负责从事件轨迹中提取重复关系、稳定因果模式和未来生成规律。

第四类是结构激活算子：
\[
\mathcal{G}_{\Theta h}:
\Theta
\rightarrow
h.
\]
它不是检索某个知识条目，而是通过
\[
W_l^0
\rightarrow
W_l^0+\Delta W_l^\Theta
\]
直接塑造当前神经动力场。

第五类是结构解释算子：
\[
\mathcal{U}_{\Theta E}:
\Theta
\rightarrow
E.
\]
它用于从长期结构生成、补全或重新解释具体情景，使结构模型可以重新下投到可检索的事件层。

第六类是快速结构修正：
\[
\mathcal{U}_{h\Theta}:
h
\rightarrow
\Theta.
\]
这一算子可以存在，但其增益必须极低：
\[
\left\|
\mathcal{U}_{h\Theta}
\right\|
\ll
\left\|
\mathcal{C}_{E\Theta}
\right\|.
\]
只有在高置信、强证据和明确结构冲突条件下，当前状态才允许直接对 $\Theta$ 进行小幅修正。

于是三相记忆真正形成
\[
h
\leftrightarrow
E
\leftrightarrow
\Theta
\]
的双向动力系统，而不是单向存储层级。

\section{三相之间的训练目标与相间一致性}

$h\rightarrow E$ 的训练目标必须保证情景向量不仅保持语义，还保留时间、因果、行动和结果信息。因此可定义
\[
\mathcal{L}_{hE}
=
\lambda_1\mathcal{L}_{reconstruct}
+
\lambda_2\mathcal{L}_{contrastive}
+
\lambda_3\mathcal{L}_{temporal}
+
\lambda_4\mathcal{L}_{causal}.
\]
其中 $\mathcal{L}_{temporal}$ 保证邻近情景的顺序关系可恢复，$\mathcal{L}_{causal}$ 则约束具有稳定因果关系的事件形成可预测结构。

$E\rightarrow h$ 的目标不应是恢复原始 Token，而应恢复
\begin{statementbox}
TaskRelevantLatentState.
\end{statementbox}
因此，其损失应围绕后续任务性能、状态重构以及情景关系保持定义，而不是简单的文本重建。

$\Theta\rightarrow h$ 则更不应被理解为知识检索。结构记忆的本质是通过改变网络生成场，使当前工作记忆在不同的长期经验背景下自然进入不同状态。因此：
\[
\boxed{
\Theta\rightarrow h
=
GeneratorShapesCurrentStateFlow.
}
\]

而 $E\rightarrow\Theta$ 应主要发生在慢速 consolidation window 中。设第 $k$ 个巩固时间窗为
\[
[t_k,t_{k+1}],
\]
在该窗口内收集一组相关情景
\[
\mathcal{B}_E,
\]
随后执行
\[
\Theta_{k+1}
=
\Theta_k
-
\eta_\Theta
\nabla_\Theta
\mathcal{L}_\Theta
\left(
\mathcal{B}_E
\right).
\]
工程上可将其安排在 idle、sleep 或 checkpoint 阶段，从而真正实现三相记忆的时间尺度分离。

类似地，$\Psi$ 更新应更加谨慎。可以定义自我更新门控：
\[
\Gamma_\Psi
=
Consistency
\times
Persistence
\times
SelfRelevance.
\]
只有当
\[
\Gamma_\Psi
>
\theta_\Psi
\]
并持续满足
\[
T
>
T_\Psi
\]
时，才允许 $s_\Psi$ 发生明显更新。这样，一次偶然失败、一次成功、一次短期奖励或惩罚都无法直接改写主体身份。

\section{信息几何视角下的三相与自我}

为了进一步说明三相之间并不是简单缓存层级，可以从信息几何角度重新描述四个对象。

工作记忆 $h(t)$ 是表示空间上的即时概率分布，因此
\[
h(t)
\in
\mathcal{P}(\mathcal{Z}),
\]
其中 $\mathcal{P}(\mathcal{Z})$ 表示定义在潜在状态空间 $\mathcal{Z}$ 上的概率测度空间。

情景记忆则不是概率分布，而是有限容量的带关系图：
\[
E(t)
\in
\mathcal{G}_{M_{\max}}(\mathcal{Z}).
\]
它保存离散事件节点及时间、因果、语义关系。

结构记忆 $\Theta$ 位于生成模型的参数流形：
\[
\Theta(t)
\in
\mathcal{M}_\Theta.
\]
它代表的是能够生成当前与未来状态的慢结构。

自我 $\Psi$ 则不能简单放入这三个空间中的任何一个，而更接近于改变这些空间局部坐标与距离关系的慢结构：
\[
g^\Psi.
\]
因此，自我可以被理解为
\begin{statementbox}
改变 Agent 如何度量自身世界的慢信息几何结构.
\end{statementbox}
这为此前“自我在 h 中表现为背景引力”的直觉提供了更严格的数学版本，也解释了为何 $\Psi$ 不应该被实现成普通第四级记忆。

\section{为什么三相记忆不是三级缓存}

传统计算机缓存通常满足
\[
L_1
\leftrightarrow
L_2
\leftrightarrow
L_3,
\]
不同层之间主要区别在访问速度和容量。而三相记忆不是这种关系，因为三者的数学对象本身不同：
\[
\boxed{
h:
Distribution,
}
\]
\[
\boxed{
E:
FiniteTemporalCausalGraph,
}
\]
\[
\boxed{
\Theta:
GenerativeParameterField.
}
\]
因此三相记忆真正描述的是
\[
\boxed{
State
\rightarrow
Trajectory
\rightarrow
Generator,
}
\]
而不是
\[
FastStorage
\rightarrow
MediumStorage
\rightarrow
SlowStorage.
\]

同样，$\Psi$ 也不是第四级缓存。它既不是当前状态，也不是事件轨迹，也不是普通世界生成模型，而是
\begin{statementbox}
CrossPhaseSelfCondition.
\end{statementbox}
它改变 h 中哪些状态更容易形成，E 中哪些经历更容易保存，以及 $\Theta$ 中哪些结构更容易稳定。因此，更准确的描述是：
\[
\boxed{
\Psi
=
\text{三相记忆共同的慢边界条件与内部坐标原点}.
}
\]

\section{工程最小实现}

以上理论并不要求一开始就修改整个基础模型。第一版工程可以采用如下最小结构。

基础模型采用冻结的预训练 LLM：
\[
W^0.
\]

工作记忆直接读取基础模型各层隐藏状态：
\[
H_1,\ldots,H_L,
\]
并增加一个 State Pooler，用于计算
\[
z_h
=
\mathcal{C}_{hE}(h).
\]

情景记忆实现为固定容量 Episodic Graph Memory：
\[
|V_E|
\leq
M_{\max}.
\]
节点保存 latent vector、timestamp、importance、self relevance、decay rate、access count 与 merge mass；向量检索可以使用 HNSW 或类似近似近邻结构，而时间边、因果边和语义边作为独立 adjacency 保存。

结构记忆 $\Theta$ 可首先实现为
\[
LoRA/AdapterBank+Router,
\]
即冻结基础模型 $W^0$，仅慢速训练结构 Adapter。

自我 $\Psi$ 则实现为一个低维慢状态
\[
s_\Psi
\]
以及三个投影组件：Layer Modulator、Episodic Filter 和 Structural Regularizer。于是，一个最小实现可以表示为
\[
\boxed{
TransformerHiddenStateDistribution
+
FixedCapacityEpisodicVectorGraph
+
SlowStructuralAdapterLLM
+
CrossPhaseSelfConditioner.
}
\]

一次完整认知循环可写为：
\[
q_t
=
Q
\left(
h_{t-1},\Psi_t
\right),
\]
\[
C_E(t)
=
Retrieve
\left(
E_t,q_t,\Psi_t
\right),
\]
\[
F_t
=
F
\left(
\Theta_t,C_E(t),\Psi_t
\right),
\]
\[
\partial_s\rho_h
=
-\nabla\cdot
\left(
\rho_hF_t
\right)
+
D\nabla^2\rho_h,
\]
从而得到新的
\[
h_t.
\]
随后压缩：
\[
z_t
=
\mathcal{C}_{hE}
\left(
h_t
\right),
\]
写入或合并：
\[
E_{t+1}
=
Update_E
\left(
E_t,z_t,\Psi_t
\right).
\]
在更慢的巩固周期：
\[
\Theta_{t+\Delta_\Theta}
=
Consolidate
\left(
E,\Theta,\Psi
\right),
\]
在更慢的身份周期：
\[
\Psi_{t+\Delta_\Psi}
=
SelfUpdate
\left(
E,\Theta,\Psi
\right).
\]

因此，整个系统最终形成四层递归：
\[
h_{t+1}
=
F_h
\left(
h_t,E_t,\Theta_t,\Psi_t,x_t
\right),
\]
\[
E_{t+1}
=
F_E
\left(
E_t,h_{t+1},\Psi_t
\right),
\]
\[
\Theta_{t+1}
=
F_\Theta
\left(
\Theta_t,E_{[t-T,t]},\Psi_t
\right),
\]
\[
\Psi_{t+1}
=
F_\Psi
\left(
\Psi_t,E_{long},\Theta_t
\right).
\]

正向沉降方向为
\[
\boxed{
FastState
\rightarrow
Episode
\rightarrow
Structure
\rightarrow
Self,
}
\]
反向约束方向为
\[
\boxed{
Self
\rightarrow
Structure
\rightarrow
EpisodeSelection
\rightarrow
CurrentState.
}
\]

\section{稳定性条件与身份连续性}

本系统是否能够真正成为持续智能的记忆基础，取决于若干必须满足的稳定条件。

第一，瞬时噪声不得直接改写慢结构：
\[
\boxed{
FastNoise
\nrightarrow
SlowStructure.
}
\]

第二，重复经验应逐渐推动慢结构改变：
\[
\boxed{
RepeatedExperience
\rightarrow
SlowLearning.
}
\]

第三，自我可以改变注意、重要度和结构稳定性，但不得覆盖现实反馈：
\[
\boxed{
SelfBias
\neq
RealityOverride.
}
\]

第四，$E$ 必须保持固定容量：
\[
\boxed{
|V_E|
\leq
M_{\max},
}
\]
并通过 decay、merge、eviction 和 consolidation 保持可管理性。

第五，基础表示空间必须保持相对稳定，否则历史情景向量将与当前 h 的表示空间失去可比较性。因此第一代系统中冻结 $W^0$ 并仅慢速训练 Adapter 是一种重要稳定策略。

第六，也是最核心的一条，是时间尺度必须真正分离：
\[
\boxed{
\tau_h
\ll
\tau_E
\ll
\tau_\Theta
\ll
\tau_\Psi.
}
\]
如果所有变量都以相似学习率端到端更新，则系统会发生 timescale collapse：当前任务损失会同时拖动 E、$\Theta$ 和 $\Psi$，三相记忆失去相态区分，身份连续性也随之失效。

因此，身份连续性不要求
\[
\Psi_t
=
\Psi_{t+\Delta}
\]
严格不变，而要求其漂移满足
\[
\left\|
\dot\Psi
\right\|
\ll
\left\|
\dot\Theta
\right\|
\ll
\left\|
\dot E
\right\|
\ll
\left\|
\dot h
\right\|.
\]
在有限正常时间窗内，还应满足
\[
D
\left(
\Psi_t,\Psi_{t+\Delta}
\right)
<
\varepsilon_\Psi.
\]
因此，长期身份更准确的数学定义是
\[
\boxed{
IdentityContinuity
=
BoundedSlowDrift,
}
\]
而不是 Static Identity。

\section{从三相记忆概念到多时间尺度神经动力系统}

本章完成了三相记忆与自我从概念结构向纯神经向量动力学实现的第一次完整转换。其最重要的结果，是明确了四个对象并非四种同类存储，而是四种不同数学结构。

工作记忆
\[
\boxed{
h(t)
=
\text{LLM 当前全部层级激活在表示空间中的状态分布}
}
\]
回答的是
\begin{statementbox}
WhatIsActiveNow?
\end{statementbox}

情景记忆
\[
\boxed{
E(t)
=
\text{固定容量的时间--因果多关系情景向量图}
}
\]
回答的是
\begin{statementbox}
WhatHappenedAlongMyTrajectory?
\end{statementbox}

结构记忆
\[
\boxed{
\Theta(t)
=
\text{由长期情景巩固形成的慢速结构生成 LLM}
}
\]
回答的是
\[
\boxed{
WhatUsuallyGeneratesWhat,
AndWhatMayComeNext?
}
\]

而自我
\[
\boxed{
\Psi(t)
}
\]
不是第四级记忆，而是同时在三个相态中发挥不同作用的慢变量。在 h 中，它表现为
\[
\boxed{
PotentialField;
}
\]
在 E 中，它表现为
\[
\boxed{
ExperienceSelectionField;
}
\]
在 $\Theta$ 中，它表现为
\begin{statementbox}
StructuralCoordinateOrigin.
\end{statementbox}

因此，自我的严格数学含义可以重新表述为：

\begin{statementbox}
Self 不是被记住的一个对象，而是当前神经状态如何形成、经历如何获得长期意义、结构记忆如何围绕同一主体组织起来的共同慢变量。
\end{statementbox}

整个三相--自我系统最终可以压缩成
\[
\boxed{
\mathcal{A}(t)
=
\left(
\rho_h,
\mathcal{G}_E,
\Theta,
\Psi
\right),
}
\]
并满足
\[
\partial_s\rho_h
=
\mathcal{F}_h
\left(
\rho_h,\mathcal{G}_E,\Theta,\Psi
\right),
\]
\[
\dot{\mathcal{G}}_E
=
\epsilon_E
\mathcal{F}_E
\left(
\mathcal{G}_E,\rho_h,\Psi
\right),
\]
\[
\dot\Theta
=
\epsilon_\Theta
\mathcal{F}_\Theta
\left(
\Theta,\mathcal{G}_E,\Psi
\right),
\]
\[
\dot\Psi
=
\epsilon_\Psi
\mathcal{F}_\Psi
\left(
\Psi,\mathcal{G}_E,\Theta
\right),
\]
其中
\[
1
\gg
\epsilon_E
\gg
\epsilon_\Theta
\gg
\epsilon_\Psi
>
0.
\]

由此，三相记忆的真正动力学主线可以写成：
\begin{statementbox}
瞬时神经状态形成情景，情景在时间中形成轨迹，轨迹经过巩固形成生成结构，生成结构重新塑造未来瞬时状态，而自我作为跨三相的慢势场、经历筛选器与结构原点，使这一持续变化始终保持为同一个主体的生命周期。
\end{statementbox}

本章至此仍只解决“记忆怎样自然演化”的问题，而尚未引入“智能体怎样主动观察和调节这一演化”。因此，后续章节可以在本章建立的
\[
h/E/\Theta/\Psi
\]
连续动力系统之上进一步引入 SPC 与 TCE，使智能体开始主动估计自身认知状态，并主动控制三相之间的召回、写入、巩固与结构迁移。那将标志着理论从“持续记忆动力学”正式进入“主动认知控制动力学”。
